% GENERATED FILE. Edit canonical lessons, then run npm run generate:books.
% canonical-source-hash: fnv1a64:694786cb9c6aba65
% canonical-lessons: JA-C500-omocha, JA-C500-nawatobi, JA-C500-kakekko, JA-C500-nazonazo, JA-C500-suberidai

\chapter{Toys and Games}
\label{ch:ja-toys-and-games}
\hlchaptermodality{500}

\begin{chapteropening}
\textbf{By the end of this chapter:} \emph{I can talk about children's games.}

Complete the last lesson of chapter 500: I can talk about children's games.
\end{chapteropening}

\section[\texorpdfstring{\ja{おもちゃ} (omocha)}{おもちゃ (omocha)}]{\ja{おもちゃ} (omocha) — a toy}

\label{lesson:JA-C500-omocha}



\begin{quote}
Before the new one: say the Japanese for a school outing, then the Japanese for a pleasure trip.
\end{quote}

\subsection*{You'll want to know: \ja{おもちゃ}}
\textbf{\ja{おもちゃ}} — \emph{omocha} — \textquotedblleft{}a toy\textquotedblright{}.

\begin{quote}
\textbf{\ja{こどもは} \ja{おもちゃで} \ja{あそびます。}} — \emph{kodomo wa omocha de asobimasu} — \textquotedblleft{}Children play with toys.\textquotedblright{}
\end{quote}

\begin{cousinweb}
\textbf{\ja{おもちゃ}} (\emph{omocha}) is a native Japanese word (wago), inherited rather than borrowed from Chinese, and usually written in kana; its older origin is not settled (JMdict).
\end{cousinweb}

\begin{grammarlens}[title={What you've built}]
One more word for describing what you lived and felt.
\end{grammarlens}

\subsection*{Your turn}
\begin{itemize}
\raggedright
  \item \emph{Say it:} \emph{omocha}
  \item \emph{Say it:} \emph{omocha}, once more
  \item \emph{Say it:} \emph{kodomo wa omocha de asobimasu} — Children play with toys.
  \item \emph{Recall:} say the Japanese for a school outing, then the Japanese for a pleasure trip, then say \emph{omocha} again
\end{itemize}

\begin{tcolorbox}[breakable,colback=teal!4,colframe=teal!35!black,title={Before you move on}]
What does \ja{おもちゃ} mean? (\textquotedblleft{}A toy\textquotedblright{}.) Say it once more.
\end{tcolorbox}
\section[\texorpdfstring{\ja{なわとび} (nawatobi)}{なわとび (nawatobi)}]{\ja{なわとび} (nawatobi) — skipping rope}

\label{lesson:JA-C500-nawatobi}



\begin{quote}
Before the new one: say the Japanese for a pleasure trip, then the Japanese for a toy.
\end{quote}

\subsection*{You'll want to know: \ja{なわとび}}
\textbf{\ja{なわとび}} — \emph{nawatobi} — \textquotedblleft{}skipping rope\textquotedblright{}.

\begin{quote}
\textbf{\ja{こうえんで} \ja{なわとびを} \ja{しました。}} — \emph{kōen de nawatobi o shimashita} — \textquotedblleft{}I skipped rope in the park.\textquotedblright{}
\end{quote}

\begin{cousinweb}
\textbf{\ja{なわとび}} (\emph{nawatobi}) is a native Japanese compound (wago), inherited rather than borrowed from Chinese: built from \emph{nawa}, \textquotedblleft{}rope\textquotedblright{} and \emph{tobi}, \textquotedblleft{}jumping\textquotedblright{} (JMdict, KANJIDIC2).
\end{cousinweb}

\begin{grammarlens}[title={What you've built}]
One more word for describing what you lived and felt.
\end{grammarlens}

\subsection*{Your turn}
\begin{itemize}
\raggedright
  \item \emph{Say it:} \emph{nawatobi}
  \item \emph{Say it:} \emph{nawatobi}, once more
  \item \emph{Say it:} \emph{kōen de nawatobi o shimashita} — I skipped rope in the park.
  \item \emph{Recall:} say the Japanese for a pleasure trip, then the Japanese for a toy, then say \emph{nawatobi} again
\end{itemize}

\begin{tcolorbox}[breakable,colback=teal!4,colframe=teal!35!black,title={Before you move on}]
What does \ja{なわとび} mean? (\textquotedblleft{}Skipping rope\textquotedblright{}.) Say it once more.
\end{tcolorbox}
\section[\texorpdfstring{\ja{かけっこ} (kakekko)}{かけっこ (kakekko)}]{\ja{かけっこ} (kakekko) — a running race (between children)}

\label{lesson:JA-C500-kakekko}



\begin{quote}
Before the new one: say the Japanese for a toy, then the Japanese for skipping rope.
\end{quote}

\subsection*{You'll want to know: \ja{かけっこ}}
\textbf{\ja{かけっこ}} — \emph{kakekko} — \textquotedblleft{}a running race (between children)\textquotedblright{}.

\begin{quote}
\textbf{\ja{こどもが} \ja{こうえんで} \ja{かけっこを} \ja{しています。}} — \emph{kodomo ga kōen de kakekko o shiteimasu} — \textquotedblleft{}Children are having a race in the park.\textquotedblright{}
\end{quote}

\begin{cousinweb}
\textbf{\ja{かけっこ}} (\emph{kakekko}) is a native Japanese word (wago), inherited rather than borrowed from Chinese, and usually written in kana; its older origin is not settled (JMdict).
\end{cousinweb}

\begin{grammarlens}[title={What you've built}]
One more word for describing what you lived and felt.
\end{grammarlens}

\subsection*{Your turn}
\begin{itemize}
\raggedright
  \item \emph{Say it:} \emph{kakekko}
  \item \emph{Say it:} \emph{kakekko}, once more
  \item \emph{Say it:} \emph{kodomo ga kōen de kakekko o shiteimasu} — Children are having a race in the park.
  \item \emph{Recall:} say the Japanese for a toy, then the Japanese for skipping rope, then say \emph{kakekko} again
\end{itemize}

\begin{tcolorbox}[breakable,colback=teal!4,colframe=teal!35!black,title={Before you move on}]
What does \ja{かけっこ} mean? (\textquotedblleft{}A running race (between children)\textquotedblright{}.) Say it once more.
\end{tcolorbox}
\section[\texorpdfstring{\ja{なぞなぞ} (nazonazo)}{なぞなぞ (nazonazo)}]{\ja{なぞなぞ} (nazonazo) — a riddle}

\label{lesson:JA-C500-nazonazo}



\begin{quote}
Before the new one: say the Japanese for skipping rope, then the Japanese for a running race (between children).
\end{quote}

\subsection*{You'll want to know: \ja{なぞなぞ}}
\textbf{\ja{なぞなぞ}} — \emph{nazonazo} — \textquotedblleft{}a riddle\textquotedblright{}.

\begin{quote}
\textbf{\ja{そふは} \ja{なぞなぞが} \ja{すきです。}} — \emph{sofu wa nazonazo ga suki desu} — \textquotedblleft{}My grandfather loves riddles.\textquotedblright{}
\end{quote}

\begin{cousinweb}
\textbf{\ja{なぞなぞ}} (\emph{nazonazo}) is a native Japanese word (wago), inherited rather than borrowed from Chinese, and usually written in kana; its older origin is not settled (JMdict).
\end{cousinweb}

\begin{grammarlens}[title={What you've built}]
One more word for describing what you lived and felt.
\end{grammarlens}

\subsection*{Your turn}
\begin{itemize}
\raggedright
  \item \emph{Say it:} \emph{nazonazo}
  \item \emph{Say it:} \emph{nazonazo}, once more
  \item \emph{Say it:} \emph{sofu wa nazonazo ga suki desu} — My grandfather loves riddles.
  \item \emph{Recall:} say the Japanese for skipping rope, then the Japanese for a running race (between children), then say \emph{nazonazo} again
\end{itemize}

\begin{tcolorbox}[breakable,colback=teal!4,colframe=teal!35!black,title={Before you move on}]
What does \ja{なぞなぞ} mean? (\textquotedblleft{}A riddle\textquotedblright{}.) Say it once more.
\end{tcolorbox}
\section[\texorpdfstring{\ja{すべりだい} (suberidai)}{すべりだい (suberidai)}]{\ja{すべりだい} (suberidai) — a slide (in a playground)}

\label{lesson:JA-C500-suberidai}



\begin{quote}
Before the new one: say the Japanese for a running race (between children), then the Japanese for a riddle.
\end{quote}

\subsection*{You'll want to know: \ja{すべりだい}}
\textbf{\ja{すべりだい}} — \emph{suberidai} — \textquotedblleft{}a slide (in a playground)\textquotedblright{}.

\begin{quote}
\textbf{\ja{こうえんの} \ja{すべりだいは} \ja{たかいです。}} — \emph{kōen no suberidai wa takai desu} — \textquotedblleft{}The slide in the park is high.\textquotedblright{}
\end{quote}

\begin{cousinweb}
\textbf{\ja{すべりだい}} (\emph{suberidai}) mixes the two layers, native and Sino-Japanese: native \emph{suberi}, \textquotedblleft{}sliding\textquotedblright{}, from \emph{suberu}, \textquotedblleft{}to slide\textquotedblright{}, and Sino-Japanese \emph{dai}, \textquotedblleft{}a stand\textquotedblright{}, borrowed from Chinese (JMdict, KANJIDIC2).
\end{cousinweb}

\begin{grammarlens}[title={What you've built}]
One more word for describing what you lived and felt.
\end{grammarlens}

\subsection*{Your turn}
\begin{itemize}
\raggedright
  \item \emph{Say it:} \emph{suberidai}
  \item \emph{Say it:} \emph{suberidai}, once more
  \item \emph{Say it:} \emph{kōen no suberidai wa takai desu} — The slide in the park is high.
  \item \emph{Recall:} say the Japanese for a running race (between children), then the Japanese for a riddle, then say \emph{suberidai} again
\end{itemize}

\begin{tcolorbox}[breakable,colback=teal!4,colframe=teal!35!black,title={Before you move on}]
What does \ja{すべりだい} mean? (\textquotedblleft{}A slide (in a playground)\textquotedblright{}.) Say it once more.
\end{tcolorbox}
